%% ======================================================================
%%  《软件学报》LaTeX 模板 —— 示例文件
%%  编译：XeLaTeX（Overleaf: Menu → Compiler → XeLaTeX）
%%  想用与 Word 同款的系统中文字体（宋体/黑体/楷体/仿宋），按系统选择：
%%      Windows: \documentclass[fontset=windows]{jos}   % 宋体 SimSun / 黑体 SimHei / 楷体 KaiTi / 仿宋 FangSong
%%      macOS:   \documentclass[fontset=mac]{jos}       % 华文宋体 / 黑体 / 华文楷体 / 华文仿宋
%%  不加选项时，Overleaf 等 Linux 环境默认用 Fandol 字体（FandolKai 也是楷体）。
%%  注：《软件学报》中文正文使用半角标点“,”“.”“:”“;”与“()”，示例遵循此约定。
%% ======================================================================

\documentclass{jos}

%% ---------- 页眉与期刊信息（录用后由编辑部提供，投稿时可保持默认）----------
\startpage{2204}                                  % 起始页码
\journalinfo{Vol.32, No.7, July 2021}             % 偶数页页眉中的卷期
\doi{10.13328/j.cnki.jos.000000}
\headauthor{张三~等}                              % 奇数页页眉：作者 等
\headtitle{论文中文题目}                          % 奇数页页眉：题目

%% ---------- 中文信息 ----------
\title{论文中文题目: 不超过 20 个汉字}
\author{张三\affmark{1},\enspace 李四\affmark{1},\enspace 王五\affmark{2},\enspace 赵六\affmark{3}}
\affiliation{1}{某某大学\enspace 计算机科学与工程学院,江苏\enspace 南京\quad 210094}
\affiliation{2}{某某大学\enspace 航天学院,江苏\enspace 南京\quad 210016}
\affiliation{3}{某某大学\enspace 计算机科学与技术学院,江苏\enspace 南京\quad 211106}
\corresponding{王五}{wangwu@example.edu.cn}

\cnabstract{摘要应包括研究的目的、方法、结果和结论,一般为 200$\sim$300 字.
摘要内容使用 9pt 楷体,“摘\ 要”标签使用黑体.这是一段用于展示排版效果的示例文字,
可直接替换为论文实际内容.这是一段用于展示排版效果的示例文字,可直接替换为论文实际内容.}
\cnkeywords{关键词一;关键词二;关键词三}
\clc{TP311}

\cncitation{张三,李四,王五,赵六.论文中文题目.软件学报,2021,32(7):2204$-$2221.
  http://www.jos.org.cn/1000-9825/0000.htm}
\encitation{Zhang S, Li S, Wang W, Zhao L. English title of the paper.
  Ruan Jian Xue Bao/Journal of Software, 2021, 32(7): 2204$-$2221 (in Chinese).
  http://www.jos.org.cn/1000-9825/0000.htm}

%% ---------- 英文信息 ----------
\entitle{English Title of the Paper}
\enauthor{ZHANG San\affmark{1},\enspace LI Si\affmark{1},\enspace WANG Wu\affmark{2},\enspace ZHAO Liu\affmark{3}}
\enaffiliation{1}{School of Computer Science and Engineering, XX University, Nanjing 210094, China}
\enaffiliation{2}{College of Aerospace Engineering, XX University, Nanjing 210016, China}
\enaffiliation{3}{College of Computer Science and Technology, XX University, Nanjing 211106, China}
\enabstract{The English abstract should correspond to the Chinese abstract. It is typeset in
7.5pt Times with 11.7pt line spacing. This is placeholder text that shows the layout of the
abstract block, and it can be replaced by the actual content of the paper. This is placeholder
text that shows the layout of the abstract block.}
\enkeywords{keyword one; keyword two; keyword three}

%% ---------- 首页脚注 ----------
\fund{国家自然科学基金 (61906090, 61773208); 江苏省自然科学基金 (BK20191287)}
\dates{2020-04-13}{2020-10-26}{2020-12-14}

\begin{document}
\maketitle

%% 《软件学报》引言不加标题、不编号,正文从第 1 节开始编号
引言部分直接书写,不加“引言”标题.正文使用 9pt 宋体与 Times New Roman,行距 14.25pt,
首行缩进 2 字.引用文献使用上标形式\cite{ref1},多篇文献会自动合并\cite{ref1,ref2,ref3}.
若需非上标形式,可写作“文献\inlinecite{ref2}指出……”.

图\ref{fig:framework}展示了插图的写法:图题位于图下方,使用 9pt 宋体,居中.
这是一段用于展示排版效果的示例文字,可直接替换为论文实际内容.这是一段用于展示排版效果的示例文字,
可直接替换为论文实际内容.

\begin{figure}[htbp]
  \centering
  \includegraphics[width=0.7\linewidth,height=4cm]{example-image}
  \caption{方法框架图}
  \label{fig:framework}
\end{figure}

\section{相关工作}
一级标题为 10.5pt 黑体,编号加粗.这是一段用于展示排版效果的示例文字,可直接替换为论文实际内容.
这是一段用于展示排版效果的示例文字,可直接替换为论文实际内容.

\subsection{二级标题}
二级标题为 9pt 黑体.这是一段用于展示排版效果的示例文字,可直接替换为论文实际内容.

\subsubsection{三级标题}
三级标题为 9pt 宋体.行内公式如 $p(\boldsymbol{z}\mid\boldsymbol{x})$,行间公式如式(\ref{eq:bayes}):
\begin{equation}
  p(\boldsymbol{z}\mid\boldsymbol{x})=\frac{p(\boldsymbol{x}\mid\boldsymbol{z})\,p(\boldsymbol{z})}{p(\boldsymbol{x})}.
  \label{eq:bayes}
\end{equation}

\section{方法}
\begin{definition}[示例定义]
定义、定理等环境的标签使用黑体,内容使用宋体.
\end{definition}

\begin{theorem}
这是一个定理的示例.
\end{theorem}
\begin{proof}
这是证明的示例.
\end{proof}

列表采用“(1)”编号:
\begin{enumerate}
  \item 第一步,这是一段用于展示排版效果的示例文字;
  \item 第二步,这是一段用于展示排版效果的示例文字.
\end{enumerate}

\section{实验分析}
表\ref{tab:data}为三线表示例:表题位于表上方,黑体 9pt,编号加粗;表内文字 7.5pt.

\begin{table}[htbp]
  \centering
  \caption{实验数据集}
  \label{tab:data}
  \begin{tabular}{ccccc}
    \toprule
    分组 & 项目 & 度量数 & 模块数 & 缺陷比例(\%) \\
    \midrule
    \multirow{2}{*}{NASA} & CM1 & 37 & 327 & 12.84 \\
                          & MW1 & 37 & 253 & 10.67 \\
    SOFTLAB & AR1 & 29 & 121 & 7.44 \\
    AEEEM   & EQ  & 61 & 324 & 39.81 \\
    \bottomrule
  \end{tabular}
\end{table}

图\ref{fig:sub}为子图示例,子图题为 7.5pt.

\begin{figure}[htbp]
  \centering
  \begin{subfigure}{0.45\linewidth}
    \centering
    \includegraphics[width=\linewidth,height=3.2cm]{example-image-a}
    \caption{PDE$\Rightarrow$PC4}
  \end{subfigure}\hfill
  \begin{subfigure}{0.45\linewidth}
    \centering
    \includegraphics[width=\linewidth,height=3.2cm]{example-image-b}
    \caption{AR6$\Rightarrow$PC4}
  \end{subfigure}
  \caption{不同权重下的性能变化}
  \label{fig:sub}
\end{figure}

\section{总\hspace{1em}结}
这是一段用于展示排版效果的示例文字,可直接替换为论文实际内容.

%% ---------- 参考文献：英文 References ----------
%% 也可改用 BibTeX: \bibliographystyle{unsrtnat}\bibliography{refs}（版式相同,条目格式取决于 bst）
\begin{thebibliography}{99}
\bibitem{ref1} Hall T, Beecham S, Bowes D, Gray D, Counsell S. A systematic literature review on fault
  prediction performance in software engineering. IEEE Trans. on Software Engineering, 2012, 38(6): 1276$-$1304.
\bibitem{ref2} Chen X, Gu Q, Liu WS, Liu SL, Ni C. Survey of static software defect prediction. Ruan Jian
  Xue Bao/Journal of Software, 2016, 27(1): 1$-$25 (in Chinese with English abstract).
  http://www.jos.org.cn/1000-9825/4923.htm
\bibitem{ref3} Kingma DP, Welling M. Auto-encoding variational Bayes. In: Proc. of the 2nd Int'l Conf.
  on Learning Representations. 2014.
\end{thebibliography}

%% ---------- 附中文参考文献（编号与上面一致）----------
\begin{cnreferences}
\cnbibitem{2} 陈翔,顾庆,刘望舒,刘树龙,倪超.静态软件缺陷预测方法研究.软件学报,2016,27(1):1$-$25.
  http://www.jos.org.cn/1000-9825/4923.htm
\end{cnreferences}

%% ---------- 作者简介 ----------
\begin{biographies}
\bio{example-image-a}{张三(1983$-$),男,博士,副教授,CCF 高级会员,主要研究领域为机器学习,数据挖掘.}
\bio{example-image-b}{李四(1995$-$),男,硕士生,主要研究领域为软件缺陷预测.}
\bio{example-image-c}{王五(1981$-$),女,博士,副研究员,CCF 专业会员,主要研究领域为机器学习,软件安全性.}
\bio{example-image}{赵六(1965$-$),男,博士,教授,博士生导师,CCF 杰出会员,主要研究领域为软件工程.}
\end{biographies}

\end{document}
